\section{CAP, PACELC y comparación}

\begin{frame}{CAP/PACELC: no etiquetamos a Netflix completo}
\begin{block}{Caso delimitado: EVCache para datos no críticos}
\textbf{Partición:} operación local disponible; replicación remota asíncrona.
\quad $\Rightarrow$ aproximación AP para este caso.

\textbf{Sin partición:} lectura local de baja latencia; posible diferencia
temporal entre regiones. \quad $\Rightarrow$ aproximación EL.
\end{block}
\begin{alertblock}{Límite}
gRPC no decide CAP. Otros almacenes de Netflix usan cuórum y garantías de
lectura después de escritura.
\end{alertblock}
{\footnotesize Fuentes: Netflix Technology Blog \cite{netflix-evcache,netflix-media-db}.}
\end{frame}

\begin{frame}{Netflix frente a infraestructura de video de YouTube}
\small
\begin{tabularx}{\linewidth}{@{}p{2.1cm}XX@{}}
\toprule
 & \textbf{Netflix} & \textbf{YouTube/Google} \\
\midrule
Arquitectura & Microservicios, cachés y almacenes distribuidos & Software
distribuido y aceleradores VCU \\
Comunicación & gRPC documentado entre backends & La fuente no identifica un RPC
concreto \\
Trade-off & Latencia frente a consistencia en EVCache & Eficiencia frente a
especialización de hardware \\
\bottomrule
\end{tabularx}
\medskip
Google reportó entre \textbf{20 y 33 veces} mayor eficiencia frente a su
plataforma previa \cite{google-vcu}.
\end{frame}
